\section{3. Requirements}

Each requirement has an ID, a priority (M~must, S~should, C~could), a criterion
that can be measured, and the tests of Section 10 that verify it. The design
decisions in Section 4.4 cite these IDs, and Section 10 says for each one
whether it was met. All times are model time.

\todo{The priorities are a proposal. Decide them together.}

\subsection{Functional}

\begin{longtable}[]{@{}L{0.07\linewidth}L{0.04\linewidth}L{0.66\linewidth}L{0.13\linewidth}@{}}
\toprule
\textbf{ID} & \textbf{Pri} & \textbf{Requirement and acceptance criterion} & \textbf{Verified by} \\
\midrule
\endhead
FR-1 & M & Physics simulates \COtwo{} (mass balance) and temperature (heat
balance) per room, driven by the true occupancy and by the airflow and setpoint
the actuators have reached. A step response matches the analytic time constant
within 5~\%; a clock jump of hours gives no overshoot. & T-01, T-02 \\
FR-2 & M & Each room has a \COtwo{} sensor, a temperature sensor and a people
counter. Each samples its room every 60~s and reports with noise: \COtwo{}
±(30~ppm + 3~\%), temperature ±0.3~°C, count ±1. A fault (stuck, dropout,
drift) can be injected into one sensor without touching the others. & T-03, T-04 \\
FR-3 & M & The controller sets the airflow and the heating setpoint of every
room once per control period (60~s), and the next readings reflect it: the loop
closes. & T-05 \\
FR-4 & M & The controller forecasts each room's occupancy one lead time ahead
(the room's \COtwo{} time constant, 10--30~min) from a model trained offline on
stored sensor readings. On held-out days its error is lower than that of ``same
as now''. & T-06, T-35 \\
FR-5 & M & An actuator validates every command before applying it: range,
unit, time to live, duplicates (idempotent on \code{cmd\_id}). It rejects an
invalid command with a reason and moves towards an accepted target at a fixed
rate. & T-07, T-08 \\
FR-6 & M & Safety, enforced in the damper independently of the controller:
\COtwo{} above 1100~ppm gives maximum airflow whatever the command says, and
the flow never drops below 0.35~l/s·m² in occupied hours or 0.10~l/s·m²
otherwise. & T-09, T-10, T-11 \\
FR-7 & S & Degraded mode. A sensor that is silent for more than 3 samples, or
stuck on one value, is not acted on: a bad \COtwo{} sensor means design flow, a
bad temperature sensor 21~°C, a bad counter means the forecast and \COtwo{}
only. A stuck \COtwo{} sensor is found within 5 samples, a stuck temperature
sensor within 12, a stuck counter once the room's \COtwo{} contradicts it. The
decision names the mode and the fault. & T-12, T-13, T-36 to T-38 \\
FR-8 & M & Every observation, true room state, decision and actuator event is
stored with its model time and a run id, and a run's stored records are enough
to rebuild its results. & T-14 \\
FR-9 & M & The policy is chosen by one setting: \code{constant},
\code{reactive}, \code{predictive} or \code{oracle}. Nothing else differs
between two evaluation runs. & T-15, T-32 \\
FR-10 & S & The dashboard shows, per room, measured \COtwo{}, people counted
and expected, airflow and temperature, with alerts for overrides, silent or
stuck sensors and actuators on fallback: now, in the BuildSim viewer, and over
time, up to a week back, on our own page. & T-16, T-39 \\
FR-11 & S & A device does nothing until an installer has registered its id,
kind and room. The registry rejects a room that is not in the floor plan, and
keeps BuildSim's equipment list in step with what is installed, also after a
BuildSim restart. & T-17, T-28 to T-31 \\
\bottomrule
\end{longtable}

\subsection{Non-functional}

\begin{longtable}[]{@{}L{0.07\linewidth}L{0.04\linewidth}L{0.66\linewidth}L{0.13\linewidth}@{}}
\toprule
\textbf{ID} & \textbf{Pri} & \textbf{Requirement and acceptance criterion} & \textbf{Verified by} \\
\midrule
\endhead
NFR-1 & M & \textbf{Air quality.} \COtwo{} is at or below 1000~ppm in at least
95~\% of occupied room-minutes over five simulated weekdays
(\code{predictive}). & T-18 \\
NFR-2 & M & \textbf{Comfort.} The temperature is within 20--24~°C in at least
95~\% of occupied room-minutes. & T-19 \\
NFR-3 & S & \textbf{Energy.} The energy proxy (fans, heating of supply air,
radiators) of \code{predictive} is at least 30~\% below \code{constant}, and
no higher than \code{reactive} while meeting NFR-1 better than
\code{reactive} does. & T-20 \\
NFR-4 & S & \textbf{Latency.} A command shows up in BuildSim within one
control period (60~s, which is 1~s real at factor 60); reported as the 95th
percentile. & T-21 \\
NFR-5 & S & \textbf{Scale.} With 50 rooms (150 sensor writes per sample,
which is 150 per real second at factor 60), a full sample cycle finishes
within one sample period. & T-22 \\
NFR-6 & M & \textbf{Robustness.} Any one of our containers can be killed and
restarted without an unsafe state: a damper without valid commands returns to
design flow, and the loop resumes within 30~s real. & T-13, T-17, T-23, T-24 \\
NFR-7 & M & \textbf{Reproducibility.} The same seed and date give the same
occupancy in every room at every minute, and every number in this report
traces to a run id. & T-25 \\
NFR-8 & C & \textbf{Evaluation cost.} One policy over five simulated weekdays
runs in at most 2.5~h real. & T-26 \\
\bottomrule
\end{longtable}

Where the numbers come from. 1000~ppm is the usual indoor-air guideline. The
band 20--24~°C is the heating-season comfort range of EN~16798-1 category~II.
The airflows come from the Swedish workplace rule and the building code
(Section 6). NFR-3 is relative on purpose: the proposal treated the gap between
reactive control and the oracle as a result to report, not as a target.
